\textbf{Задача.}
Пусть для некоторого фиксированного $d > 1$ выполнено $\bpp_{0,1/2^{n^d}}=\rp$. Докажите, что тогда $\rp =\np$.

\begin{proof}
Знаем, что $\rp\subset\np$ (сертификат --- это успешная реализации случайных бит)

Достаточно показать, что $\sat\in\rp$, так как $\sat \in \npc$, а $\rp$ замкнут относительно сводимости. Последнее проверяется просто: если имеется сводимость $A \le_p B$ функцией $f$, где $B \in \rp$, то вероятностный алгоритм для $A$ запускает функцию $f$ с входом в качестве аргумента, а потом -- вероятностный алгоритм для языка $B$.

$\sat \in \bpp_{0,\frac{1}{2^n}}$. Достаточно брать значения переменных как случайные и возвращать значение формулы при этих значениях. Если формула выполнима, то существует хотя бы один выполняющий набор из $2^n$ наборов, где $n$ --- число переменных в формуле.

Также $\bpp_{0,\frac{1}{2^n}} \subset \bpp_{0,\frac{1}{2^{n^d}}} \ \forall d \ge 0$. В условии сказано, что $d > 1$, но мы рассмотрим и случай, когда $0 < d \le 1$.

\begin{enumerate}
    \item Для $d \ge 1$ верно $\Pr_r[V(x, r) = 1] \ge \frac{1}{2^n} \ge \frac{1}{2^{n^d}}$. Поэтому вложение верно.
    
    \item Теперь покажем для $d \in (0,1)$. Сделаем амплификацию: запустим алгоритм $k$ раз, которое выведем позже, каждый раз с новыми случайными битами и возьмём дизъюнкцию промежуточных результатов. Отвержение очевидно как дизъюнкция нулей. Теперь нужно рассмотреть вероятность ошибки. Она будет такова:
    
    \[\left(1-\frac{1}{2^n}\right)^k \le 1 - \frac{1}{2^{n^d}}\]
    
    Значит, можно подобрать $k$. (P. S. Правда я хз, почему)
\end{enumerate}
    
Получается, $\sat \in \bpp_{0,\frac{1}{2^n}} \subset \bpp_{0,\frac{1}{2^{n^d}}} = \rp$.

\end{proof}
